% =============================================================================
% CAPÍTULO 5 — DISEÑO DETALLADO DEL FRONTEND (ethoshubF)
% Archivo maestro del capítulo. Importa las secciones desde sections/cap05/.
% =============================================================================
\chapter{Diseño Detallado del Frontend (EthosHub)}
\label{ch:frontend}
\resumenCapitulo{Este capítulo documenta la SPA: la jerarquía de componentes, los flujos de
navegación por rol, el enrutamiento y la condición de Aplicación Web Progresiva, la
estrategia de interfaz (mapas, edición de código y gráficas), la gestión de estado con
Zustand, el consumo asíncrono de servicios, la internacionalización y la calidad del código.}

Este capítulo documenta el diseño interno de la \textit{Single Page Application} (SPA)
\texttt{ethoshubF}, el cliente que consume la API REST descrita en el
capítulo~\ref{ch:backend} (pág.~\pageref{ch:backend}). Se cubren, con el detalle que exige
la norma ISO/IEC/IEEE~15289:2024: la jerarquía de componentes y la composición raíz; los
flujos de navegación por rol; el sistema de enrutamiento y su mapa de sitio; la condición
de \textbf{Aplicación Web Progresiva} (PWA) y la estrategia de interfaz; la gestión de
estado con Zustand; el ciclo de vida del consumo asíncrono de servicios; y las
integraciones avanzadas (mapas, edición de código, gráficas y asistencia de IA).

El código está escrito en \textbf{TypeScript estricto} y construido con \textbf{Vite~5}.
Se organiza bajo el paradigma \textbf{\textit{feature-first}}, con el alias de módulo
\comando{@ $\rightarrow$ ./src} que evita las rutas relativas frágiles. El \textit{stack}
de presentación combina \textbf{React~18.3}, \textbf{React Router~6}, \textbf{Zustand}
para el estado, \textbf{Axios} para el transporte HTTP, \textbf{TailwindCSS} y primitivas
\textbf{Radix UI} para la capa visual, e \textbf{i18next} para la internacionalización
(español, inglés y portugués).

\begin{NotaTecnica}
La SPA no se despliega de forma independiente: el \textit{script} de construcción compila
el \textit{bundle} de Vite a \pathbreak{dist/} y este se copia al directorio
\pathbreak{static/} del backend, que lo sirve desde el mismo \textit{JAR} monolítico. El
filtro \texttt{SpaFallbackFilter} (capítulo~\ref{ch:backend}, sección~\ref{subsec:filtros},
pág.~\pageref{subsec:filtros}) reenvía las rutas de navegación a \texttt{index.html},
permitiendo el enrutamiento del lado del cliente.
\end{NotaTecnica}

% --- Secciones del capítulo (orden del índice instrucciones.md §5) ---
% =============================================================================
% 5.1 ÁRBOL Y JERARQUÍA DE COMPONENTES DE LA APLICACIÓN
% =============================================================================
\section{Árbol y Jerarquía de Componentes de la Aplicación}
\label{sec:arbol_componentes}

La SPA se estructura como un árbol de componentes React cuya raíz establece el contexto
global y cuyas hojas son las vistas concretas por rol. Esta sección documenta esa
jerarquía, desde el punto de entrada hasta los componentes reutilizables.

La arquitectura del frontend impone una \textbf{dirección de dependencias} estricta entre
capas: las capas superiores dependen de las inferiores, nunca al revés. La
figura~\ref{fig:capas-fe} representa este flujo, que evita ciclos y mantiene navegable el
código.

\vspace{0.3cm}
\begin{figure}[h!]
\centering
\begin{tikzpicture}[
    capa/.style={rectangle, rounded corners=3pt, draw=c4Sistema!70!black, font=\sffamily\scriptsize,
        align=center, minimum width=8.5cm, minimum height=0.75cm, inner sep=3pt},
    rel/.style={-{Stealth[length=2mm]}, semithick, draw=grisISO}]
    \node[capa, fill=c4Sistema, text=white] (app) at (0,3.6) {\texttt{app/} — providers · router · layouts};
    \node[capa, fill=c4Contenedor, text=white] (pages) at (0,2.55) {\texttt{pages/} — vistas por rol};
    \node[capa, fill=c4Componente] (feat) at (0,1.5) {\texttt{features/} — lógica de dominio};
    \node[capa, fill=c4Componente!60] (shared) at (0,0.45) {\texttt{components/} · \texttt{shared/ui} — componentes};
    \node[capa, fill=c4Datos!25] (svc) at (0,-0.6) {\texttt{store/} · \texttt{shared/services} — estado y API};

    \draw[rel] (app) -- (pages); \draw[rel] (pages) -- (feat);
    \draw[rel] (feat) -- (shared); \draw[rel] (shared) -- (svc);
    \draw[rel] (pages.east) to[bend left=55] (svc.east);
\end{tikzpicture}
\caption{Dirección de dependencias entre capas del frontend.}
\label{fig:capas-fe}
\end{figure}

\subsection{Estructura raíz: \texttt{app/providers} y \texttt{app/router}}
\label{subsec:estructura_raiz}

El punto de entrada \comando{main.tsx} monta la aplicación dentro de una pila de
\textbf{\textit{providers}} globales que establecen el contexto compartido (autenticación,
internacionalización, notificaciones \textit{toast}), bajo la cual opera el
\textbf{router}. La figura~\ref{fig:arbol-raiz} muestra esta composición raíz.

\vspace{0.3cm}
\begin{figure}[h!]
\centering
\begin{tikzpicture}[
    n/.style={rectangle, rounded corners=2pt, draw=c4Sistema!70!black, fill=c4Componente!30,
        font=\sffamily\scriptsize, align=center, minimum width=3.0cm, minimum height=0.7cm, inner sep=3pt},
    root/.style={n, fill=c4Sistema, text=white},
    rel/.style={-{Stealth[length=2mm]}, semithick, draw=grisISO},
    level distance=1.0cm, node distance=0.5cm and 0.6cm]
    \node[root] (main) at (0,4) {\texttt{main.tsx}};
    \node[n] (prov) at (0,3) {\texttt{app/providers/}\\\scriptsize Auth · i18n · Toaster};
    \node[n] (router) at (0,2) {\texttt{app/router/}\\\scriptsize index.tsx (RouterProvider)};
    \node[n] (rsg) at (0,1) {\texttt{RoleScopeGuard}};
    \node[n] (layouts) at (-3.0,0) {\texttt{app/layouts/}};
    \node[n] (pages) at (0,0) {\texttt{pages/}};
    \node[n] (features) at (3.0,0) {\texttt{features/}};
    \node[n] (comp) at (0,-1) {\texttt{components/} (ui · auth · brand)};

    \draw[rel] (main) -- (prov);
    \draw[rel] (prov) -- (router);
    \draw[rel] (router) -- (rsg);
    \draw[rel] (rsg) -- (layouts);
    \draw[rel] (rsg) -- (pages);
    \draw[rel] (rsg) -- (features);
    \draw[rel] (pages) -- (comp);
    \draw[rel] (features) -- (comp);
\end{tikzpicture}
\caption{Composición raíz de la SPA: \textit{providers}, \textit{router} y guardas.}
\label{fig:arbol-raiz}
\end{figure}

\paragraph{Carga diferida por ruta (\textit{code splitting}).} Todas las páginas se
importan con \comando{React.lazy} y se renderizan dentro de un \comando{<Suspense>} con un
\textit{skeleton} de carga. Esto, combinado con la división manual de \textit{chunks} de
Vite (\texttt{vendor: react, react-dom, react-router-dom}), reduce el tamaño del paquete
inicial y acelera el primer renderizado.

\begin{lstlisting}[language=JavaScript, caption={Carga diferida de páginas por ruta (\texttt{app/router/index.tsx}).}, label={lst:lazy}]
const CVStudioPage = lazy(() => import('@/pages/dashboard/CVStudioPage'));

function S({ page: Page }: { page: React.ComponentType }) {
  return (
    <Suspense fallback={<PageLoader />}>
      <Page />
    </Suspense>
  );
}
\end{lstlisting}

\subsection{\textit{Layouts}, páginas y componentes encapsulados por \textit{feature}}
\label{subsec:layouts_paginas}

La organización del código separa con claridad las \textbf{capas de presentación}
(\textit{layouts} y \textit{pages}, por rol) de la \textbf{lógica de dominio}
(\textit{features}) y de los \textbf{componentes reutilizables}. La
tabla~\ref{tab:carpetas-fe} describe cada directorio de primer nivel.

\vspace{0.2cm}
\noindent
\renewcommand{\arraystretch}{1.3}
\fontsize{8.5}{10.2}\selectfont\sffamily
\begin{tabularx}{\textwidth}{|L{3.0cm}|Y|}
\hline
\rowcolor{headerTabla}
\sffamily\textbf{Directorio} & \sffamily\textbf{Responsabilidad} \\
\hline
\pathbreak{app/} & Composición de la app: \textit{providers}, \textit{layouts} por sección y \textit{router}. \\
\hline
\rowcolor{grisTecnico}
\pathbreak{pages/} & Vistas por rol: \texttt{auth}, \texttt{dashboard}, \texttt{recruiter}, \texttt{admin}, \texttt{public}. \\
\hline
\pathbreak{features/} & Lógica de dominio encapsulada: \texttt{portfolio}, \texttt{projects}, \texttt{recruiter}. \\
\hline
\rowcolor{grisTecnico}
\pathbreak{components/} & Componentes reutilizables: \texttt{ui}, \texttt{auth}, \texttt{brand}, \texttt{preferences}. \\
\hline
\pathbreak{shared/} & Servicios de API, tipos, utilidades, \textit{landing}, primitivas \texttt{ui}. \\
\hline
\rowcolor{grisTecnico}
\pathbreak{store/} & \textit{Stores} Zustand por dominio. \\
\hline
\pathbreak{hooks/} · \pathbreak{lib/} & \textit{Custom hooks} y utilidades de bajo nivel. \\
\hline
\rowcolor{grisTecnico}
\pathbreak{i18n/} & Configuración y traducciones (\texttt{es} / \texttt{en} / \texttt{pt}). \\
\hline
\end{tabularx}
\label{tab:carpetas-fe}

\paragraph{\textit{Layouts} por rol.} Existen cinco \textit{layouts} que componen la
estructura visual según el contexto de navegación, enumerados en la
tabla~\ref{tab:layouts}. Cada uno envuelve sus rutas hijas mediante el \comando{<Outlet/>}
de React Router.

\vspace{0.2cm}
\noindent
\renewcommand{\arraystretch}{1.3}
\fontsize{8.5}{10.2}\selectfont\sffamily
\begin{tabularx}{\textwidth}{|L{3.6cm}|Y|}
\hline
\rowcolor{headerTabla}
\sffamily\textbf{Layout} & \sffamily\textbf{Contexto} \\
\hline
\texttt{AuthLayout} & Pantallas de \textit{login}, registro y recuperación de contraseña. \\ \hline
\rowcolor{grisTecnico}
\texttt{DashboardLayout} & Espacio de trabajo del profesional (portafolio, CV, proyectos). \\ \hline
\texttt{RecruiterLayout} & Espacio del reclutador (descubrimiento, \textit{rankings}, chat). \\ \hline
\rowcolor{grisTecnico}
\texttt{AdminLayout} & Panel de administración (métricas, moderación). \\ \hline
\texttt{PublicPortfolioLayout} & Vista pública de portafolios (\pathbreak{/p/:slug}). \\ \hline
\end{tabularx}
\label{tab:layouts}

\paragraph{Catálogo de componentes reutilizables.} El directorio \pathbreak{components/}
agrupa los componentes por familia. La tabla~\ref{tab:componentes} cataloga las familias
y ejemplos representativos.

\vspace{0.2cm}
\noindent
\renewcommand{\arraystretch}{1.3}
\fontsize{8.5}{10.2}\selectfont\sffamily
\begin{tabularx}{\textwidth}{|L{2.8cm}|Y|}
\hline
\rowcolor{headerTabla}
\sffamily\textbf{Familia} & \sffamily\textbf{Componentes y propósito} \\
\hline
\texttt{ui} & Primitivas de interfaz sobre Radix: botones, diálogos, menús, \textit{tooltips}, \textit{skeletons}. \\ \hline
\rowcolor{grisTecnico}
\texttt{auth} & Guardas y formularios de acceso, campos de OTP (\texttt{input-otp}). \\ \hline
\texttt{brand} & Identidad visual: logotipos, marca de agua, elementos corporativos. \\ \hline
\rowcolor{grisTecnico}
\texttt{preferences} & Selectores de idioma y tema, ajustes del usuario. \\ \hline
\end{tabularx}
\label{tab:componentes}

\paragraph{Componentes de dominio (\textit{features}).} Frente a los componentes
genéricos, las \textit{features} encapsulan lógica de negocio específica: \texttt{portfolio}
(composición del portafolio público), \texttt{projects} (gestión de proyectos y su
galería de medios) y \texttt{recruiter} (tarjetas de talento, \textit{match ring} de
afinidad y notas privadas). Esta separación entre componentes genéricos y de dominio es la
materialización del paradigma \textit{feature-first} en la capa de presentación.

\subsection{\textit{Custom hooks} y lógica reutilizable}
\label{subsec:hooks}

La lógica con estado que se comparte entre componentes se extrae a \textbf{\textit{custom
hooks}}, evitando la duplicación y manteniendo los componentes centrados en la
presentación. El ejemplo canónico es \comando{useAuthFlow}, que encapsula la orquestación
del proceso de autenticación (validación, llamada al servicio, manejo de errores,
redirección por rol) y la expone como una API declarativa a las páginas de \textit{login}
y registro. El \textit{provider} \comando{AuthProvider} complementa este patrón
inyectando el contexto de sesión en todo el árbol y escuchando el evento
\comando{auth:unauthorized} emitido por el interceptor de Axios
(sección~\ref{subsec:interceptor_resp}, pág.~\pageref{subsec:interceptor_resp}) para
disparar el cierre de sesión global.

\begin{NotaTecnica}
La regla de los \textit{hooks} (invocación incondicional en el nivel superior del
componente) la verifica estáticamente el \textit{plugin} \texttt{eslint-plugin-react-hooks}
(sección~\ref{subsec:eslint}, pág.~\pageref{subsec:eslint}), de modo que una violación se
detecta en el \textit{linting} y no en tiempo de ejecución.
\end{NotaTecnica}
   % 5.1 Árbol y jerarquía de componentes
% =============================================================================
% 5.2 USER FLOWS (FLUJOS DE NAVEGACIÓN DEL PERFIL DE USUARIO)
% =============================================================================
\section{User Flows (Flujos de Navegación del Perfil de Usuario)}
\label{sec:user_flows}

Los \textit{user flows} modelan la experiencia de navegación de los dos perfiles
principales del sistema, trazando el camino desde la entrada hasta el objetivo de valor de
cada rol. Estos flujos materializan los casos de uso definidos en el
capítulo~\ref{ch:requisitos} (pág.~\pageref{ch:requisitos}).

\subsection{Registro, configuración inicial y carga de portafolio}
\label{subsec:flow_registro}

La figura~\ref{fig:flow-onboarding} representa el embudo de incorporación de un
profesional: desde el registro hasta la publicación de su portafolio, incluyendo la
ramificación por rol que ocurre inmediatamente tras el alta.

\vspace{0.3cm}
\begin{figure}[h!]
\centering
\begin{tikzpicture}[node distance=0.8cm and 1.1cm]
    \node[flujoInicio] (start) {Registro};
    \node[flujoDecision, below=of start] (rol) {¿Rol?};
    \node[flujoProceso, below left=0.9cm and 0.2cm of rol] (prof) {Dashboard\\Profesional};
    \node[flujoProceso, below right=0.9cm and 0.2cm of rol] (recl) {Dashboard\\Reclutador};
    \node[flujoProceso, below=of prof] (perfil) {Configurar perfil\\(bio, foto, skills)};
    \node[flujoProceso, below=of perfil] (proy) {Cargar proyectos\\y experiencia};
    \node[flujoProceso, below=of proy] (ajustes) {Ajustes de\\portafolio (SEO)};
    \node[flujoInicio, fill=c4Sistema, below=of ajustes] (pub) {Portafolio público};

    \draw[flujoFlecha] (start) -- (rol);
    \draw[flujoFlecha] (rol) -- node[c4etiqueta, left]{PROFESSIONAL} (prof);
    \draw[flujoFlecha] (rol) -- node[c4etiqueta, right]{RECRUITER} (recl);
    \draw[flujoFlecha] (prof) -- (perfil);
    \draw[flujoFlecha] (perfil) -- (proy);
    \draw[flujoFlecha] (proy) -- (ajustes);
    \draw[flujoFlecha] (ajustes) -- (pub);
\end{tikzpicture}
\caption{Flujo de incorporación del profesional: del registro al portafolio público.}
\label{fig:flow-onboarding}
\end{figure}

El paso de \textbf{configurar experiencia} incorpora una sub-experiencia geolocalizada:
en \pathbreak{/dashboard/experience} el usuario sitúa cada empleo en un mapa interactivo
(Leaflet) y el sistema deriva la dirección textual por geocodificación inversa, como se
detalla en la sección~\ref{subsec:comp_avanzados} (pág.~\pageref{subsec:comp_avanzados}).

\subsection{Embudo del reclutador (búsqueda, filtrado y analíticas)}
\label{subsec:flow_reclutador}

El reclutador opera un \textbf{embudo de descubrimiento}: parte de una búsqueda amplia
(filtros estructurados o lenguaje natural traducido por IA), refina hacia un
\textit{pipeline} gestionable mediante un tablero Kanban, y culmina en analíticas e
\textit{insights} generados por IA. La figura~\ref{fig:flow-recruiter} lo ilustra.

\vspace{0.3cm}
\begin{figure}[h!]
\centering
\begin{tikzpicture}[node distance=0.75cm and 1.4cm]
    \node[flujoInicio] (start) {Talent Discovery};
    \node[flujoProceso, below=of start] (search) {Buscar (filtros / NLP)};
    \node[flujoProceso, below=of search] (kanban) {Organizar en Kanban\\(drag \& drop)};
    \node[flujoProceso, below=of kanban] (notas) {Notas privadas\\+ match ring};
    \node[flujoProceso, below=of notas] (dash) {Dashboard\\(funnel · radar · KPIs)};
    \node[flujoInicio, fill=c4Sistema, below=of dash] (ai) {Insights IA + chat};

    \draw[flujoFlecha] (start) -- (search);
    \draw[flujoFlecha] (search) -- (kanban);
    \draw[flujoFlecha] (kanban) -- (notas);
    \draw[flujoFlecha] (notas) -- (dash);
    \draw[flujoFlecha] (dash) -- (ai);
\end{tikzpicture}
\caption{Embudo de descubrimiento de talento del reclutador.}
\label{fig:flow-recruiter}
\end{figure}

\begin{AvisoNota}
El paso \textbf{Buscar (NLP)} invoca el modo \texttt{nlp\_search} del motor de IA
(capítulo~\ref{ch:backend}, tabla~\ref{tab:ia-modos}, pág.~\pageref{tab:ia-modos}), que
traduce una frase como ``desarrolladores React senior en Bolivia'' a un objeto de filtros
estructurados que el frontend aplica sobre el catálogo de talento. La IA no consulta la
base de datos: solo estructura la intención del reclutador.
\end{AvisoNota}

\subsection{Visitante anónimo de un portafolio público}
\label{subsec:flow_visitante}

El tercer flujo relevante es el del \textbf{visitante no autenticado} que consulta un
portafolio compartido (\pathbreak{/p/:slug}). Este recorrido no requiere sesión y puede
estar protegido por contraseña según los ajustes del dueño. La
figura~\ref{fig:flow-visitante} lo modela.

\vspace{0.3cm}
\begin{figure}[h!]
\centering
\begin{tikzpicture}[node distance=0.75cm and 1.1cm]
    \node[flujoInicio] (start) {Visita \pathbreak{/p/:slug}};
    \node[flujoDecision, below=of start] (pwd) {¿Protegido\\por clave?};
    \node[flujoProceso, right=2.4cm of pwd] (gate) {Pantalla de\\contraseña};
    \node[flujoDecision, below=1.0cm of pwd] (pub) {¿Publicado?};
    \node[flujoProceso, right=2.4cm of pub] (404) {No encontrado\\(404)};
    \node[flujoInicio, fill=c4Sistema, below=1.0cm of pub] (view) {Portafolio renderizado};

    \draw[flujoFlecha] (start) -- (pwd);
    \draw[flujoFlecha] (pwd) -- node[c4etiqueta, above]{sí} (gate);
    \draw[flujoFlecha] (gate) |- node[c4etiqueta, near start]{ok} (pub);
    \draw[flujoFlecha] (pwd) -- node[c4etiqueta, left]{no} (pub);
    \draw[flujoFlecha] (pub) -- node[c4etiqueta, above]{no} (404);
    \draw[flujoFlecha] (pub) -- node[c4etiqueta, left]{sí} (view);
\end{tikzpicture}
\caption{Flujo del visitante anónimo a un portafolio público.}
\label{fig:flow-visitante}
\end{figure}

Este flujo consume exclusivamente \textit{endpoints} públicos
(\pathbreak{/api/v1/portfolio/public/**}) declarados en \texttt{PUBLIC\_PATHS}
(capítulo~\ref{ch:backend}, tabla~\ref{tab:public-paths}, pág.~\pageref{tab:public-paths}),
sin token; las políticas RLS del motor garantizan que solo se exponga contenido marcado
como público y publicado.
     % 5.2 User flows
% =============================================================================
% 5.3 SITEMAP, SISTEMA DE ENRUTAMIENTO Y CONDICIÓN PWA
% =============================================================================
\section{Sitemap y Sistema de Enrutamiento}
\label{sec:sitemap}

El enrutamiento gobierna qué vista se renderiza según la URL y el nivel de acceso del
usuario. Esta sección documenta el mapa de sitio completo, las guardas de protección, las
reglas responsivas y la condición de Aplicación Web Progresiva de la SPA.

\subsection{Rutas públicas, privadas y protegidas por roles (React Router 6)}
\label{subsec:rutas}

El enrutamiento usa \textbf{React Router DOM 6.22} con un árbol de rutas anidadas por
\textit{layout}. La protección se compone como envolturas (\textit{wrappers}) de tres
niveles, descritas en la tabla~\ref{tab:guardas}.

\vspace{0.2cm}
\noindent
\renewcommand{\arraystretch}{1.3}
\fontsize{8.5}{10.2}\selectfont\sffamily
\begin{tabularx}{\textwidth}{|L{3.6cm}|Y|}
\hline
\rowcolor{headerTabla}
\sffamily\textbf{Guarda} & \sffamily\textbf{Responsabilidad} \\
\hline
\texttt{ProtectedRoute} & Exige autenticación y un rol permitido (\texttt{allowedRoles}); redirige a \pathbreak{/login} si no hay sesión válida. \\ \hline
\rowcolor{grisTecnico}
\texttt{RoleScopeGuard} & Acota la navegación del administrador a \pathbreak{/admin/*}; envuelve todo el árbol de rutas. \\ \hline
\texttt{RouteErrorBoundary} & Captura errores de renderizado por ruta y muestra una vista de error en vez de una pantalla en blanco. \\ \hline
\end{tabularx}
\label{tab:guardas}

La tabla~\ref{tab:rutas-fe} cataloga las rutas reales agrupadas por \textit{layout} y
nivel de acceso, extraídas de \pathbreak{app/router/index.tsx}.

\vspace{0.2cm}
\noindent
\renewcommand{\arraystretch}{1.25}
\fontsize{8}{9.6}\selectfont\sffamily
\begin{tabularx}{\textwidth}{|L{2.6cm}|L{4.6cm}|Y|}
\hline
\rowcolor{headerTabla}
\sffamily\textbf{Acceso} & \sffamily\textbf{Ruta} & \sffamily\textbf{Vista} \\
\hline
Pública (auth) & \pathbreak{/login} · \pathbreak{/register} · \pathbreak{/forgot-password} & Entrada al sistema. \\ \hline
\rowcolor{grisTecnico}
Pública (auth) & \pathbreak{/oauth2/callback} · \pathbreak{/oauth-success} & Retorno del flujo OAuth de Supabase. \\ \hline
Pública & \pathbreak{/explorar} · \pathbreak{/p/:slug} & Exploración y portafolio público. \\ \hline
\rowcolor{grisTecnico}
Pública & \pathbreak{/privacidad} · \pathbreak{/terminos} & Páginas legales. \\ \hline
Profesional & \pathbreak{/dashboard/portfolio} · \pathbreak{/dashboard/cv-studio} & Portafolio y editor de CV. \\ \hline
\rowcolor{grisTecnico}
Profesional & \pathbreak{/dashboard/projects} · \pathbreak{/dashboard/projects/:projectId} & Proyectos y su detalle. \\ \hline
Profesional & \pathbreak{/dashboard/experience} · \pathbreak{/dashboard/education} & Experiencia y formación. \\ \hline
\rowcolor{grisTecnico}
Profesional & \pathbreak{/dashboard/chat} · \pathbreak{/dashboard/chat/:chatId} & Chat con reclutadores. \\ \hline
Reclutador & \pathbreak{/recruiter/talent-discovery} & Descubrimiento de talento. \\ \hline
\rowcolor{grisTecnico}
Reclutador & \pathbreak{/recruiter/dashboard} · \pathbreak{/recruiter/rankings} & Analíticas y \textit{leaderboard}. \\ \hline
Reclutador & \pathbreak{/recruiter/chat} & Chat con candidatos. \\ \hline
\rowcolor{grisTecnico}
Admin & \pathbreak{/admin/*} & Métricas, perfiles, moderación, \textit{skills}, correo. \\ \hline
\end{tabularx}
\label{tab:rutas-fe}

La figura~\ref{fig:sitemap} sintetiza el mapa de sitio por nivel de acceso y la guarda
responsable de cada zona.

\vspace{0.3cm}
\begin{figure}[h!]
\centering
\begin{tikzpicture}[
    grp/.style={rectangle, rounded corners=4pt, draw, font=\sffamily\scriptsize\bfseries, inner sep=8pt},
    r/.style={font=\sffamily\scriptsize, align=left, text=grisISO!90!black}]
    \node[grp, draw=flujoInicio!70!black, fill=flujoInicio!12] (pub) at (-4.6,0) {
        \begin{tabular}{l}\textbf{\textcolor{flujoInicio!60!black}{Públicas}}\\[2pt]
        \texttt{/login · /register}\\ \texttt{/forgot-password}\\ \texttt{/oauth2/callback}\\
        \texttt{/explorar · /talent}\\ \texttt{/p/:slug}\\ \texttt{/privacidad}\end{tabular}};
    \node[grp, draw=c4Sistema!70!black, fill=c4Componente!18] (priv) at (0,0) {
        \begin{tabular}{l}\textbf{\textcolor{c4Sistema}{Privadas (auth)}}\\[2pt]
        \texttt{/dashboard/*}\\ \texttt{/recruiter/*}\\ \texttt{ProtectedRoute}\\ \scriptsize(redirige a /login\\ \scriptsize si expira el JWT)\end{tabular}};
    \node[grp, draw=colorAcento!80!black, fill=colorAcento!18] (adm) at (4.4,0) {
        \begin{tabular}{l}\textbf{\textcolor{colorAcento!50!black}{Rol ADMIN}}\\[2pt]
        \texttt{/admin/*}\\ \texttt{RoleScopeGuard}\\ \scriptsize(el admin no puede\\ \scriptsize salir de /admin)\end{tabular}};
\end{tikzpicture}
\caption{Mapa de sitio por nivel de acceso y guardas de enrutamiento.}
\label{fig:sitemap}
\end{figure}

\begin{NotaTecnica}
La guarda \comando{RoleScopeGuard} implementa una regla de negocio estricta: un perfil
administrador opera \textbf{exclusivamente} dentro de \pathbreak{/admin/*}; cualquier
intento de navegar fuera lo redirige al \textit{dashboard} de administración. El acceso
inverso (un no-admin a \pathbreak{/admin}) lo bloquea \texttt{ProtectedRoute} con su lista
\texttt{allowedRoles}. Así, la frontera de roles se aplica en el cliente como primera
barrera, y se reafirma en el servidor (capítulo~\ref{ch:backend}) como autoridad final.
\end{NotaTecnica}

\paragraph{Árbol de anidamiento de rutas.} React Router 6 compone las rutas como un
\textbf{árbol}: el \texttt{RoleScopeGuard} envuelve todo, y bajo él cada \textit{layout}
agrupa sus rutas hijas, que se renderizan en el \comando{<Outlet/>} del \textit{layout}.
La figura~\ref{fig:route-tree} representa esta jerarquía.

\vspace{0.3cm}
\begin{figure}[h!]
\centering
\begin{tikzpicture}[
    n/.style={rectangle, rounded corners=2pt, draw=c4Sistema!70!black, fill=c4Componente!30,
        font=\sffamily\scriptsize, align=center, minimum height=0.6cm, inner sep=3pt},
    guard/.style={n, fill=colorAcento!25, draw=colorAcento!80!black},
    leaf/.style={n, fill=grisTecnico, draw=grisISO!60},
    rel/.style={-{Stealth[length=1.6mm]}, semithick, draw=grisISO}]
    \node[guard] (rsg) at (0,3.2) {RoleScopeGuard (raíz)};
    \node[n] (auth) at (-4.2,2.0) {AuthLayout};
    \node[guard] (pr1) at (-1.4,2.0) {ProtectedRoute\\\scriptsize professional};
    \node[guard] (pr2) at (1.8,2.0) {ProtectedRoute\\\scriptsize recruiter};
    \node[guard] (pr3) at (4.6,2.0) {ProtectedRoute\\\scriptsize admin};
    \node[leaf] (al) at (-4.2,0.7) {login · register\\forgot-password};
    \node[leaf] (dl) at (-1.4,0.7) {DashboardLayout\\\scriptsize /dashboard/*};
    \node[leaf] (rl) at (1.8,0.7) {RecruiterLayout\\\scriptsize /recruiter/*};
    \node[leaf] (adl) at (4.6,0.7) {AdminLayout\\\scriptsize /admin/*};

    \draw[rel] (rsg) -- (auth); \draw[rel] (rsg) -- (pr1);
    \draw[rel] (rsg) -- (pr2); \draw[rel] (rsg) -- (pr3);
    \draw[rel] (auth) -- (al); \draw[rel] (pr1) -- (dl);
    \draw[rel] (pr2) -- (rl); \draw[rel] (pr3) -- (adl);
\end{tikzpicture}
\caption{Árbol de anidamiento de rutas y guardas de React Router 6.}
\label{fig:route-tree}
\end{figure}

\subsection{Reglas de visualización responsiva}
\label{subsec:responsive}

El \textit{layout} general aplica una especificación responsiva deliberada respecto al
\textbf{Sidebar} de navegación, descrita en la tabla~\ref{tab:responsive}.

\vspace{0.2cm}
\noindent
\renewcommand{\arraystretch}{1.3}
\fontsize{9}{10.8}\selectfont\sffamily
\begin{tabularx}{\textwidth}{|L{2.6cm}|Y|}
\hline
\rowcolor{headerTabla}
\sffamily\textbf{Dispositivo} & \sffamily\textbf{Comportamiento del Sidebar} \\
\hline
Escritorio & La navegación principal se integra en el \textit{layout} sin un Sidebar lateral persistente, maximizando el área de contenido. \\
\hline
\rowcolor{grisTecnico}
Móvil & El Sidebar se renderiza de forma exclusiva como navegación colapsable, adaptada a la interacción táctil. \\
\hline
\end{tabularx}
\label{tab:responsive}

\subsection{Condición de Aplicación Web Progresiva (PWA)}
\label{subsec:pwa}

EthosHub es una \textbf{Aplicación Web Progresiva}: puede instalarse en la pantalla de
inicio del dispositivo y comportarse como una aplicación nativa. La condición PWA se
sustenta en el \textit{Web App Manifest} (\pathbreak{public/manifest.webmanifest}) y en
las \textit{meta} de capacidad nativa declaradas en \pathbreak{index.html}. La
tabla~\ref{tab:pwa-manifest} documenta los campos clave del manifiesto real.

\vspace{0.2cm}
\noindent
\renewcommand{\arraystretch}{1.3}
\fontsize{8.5}{10.2}\selectfont\sffamily
\begin{tabularx}{\textwidth}{|L{3.2cm}|L{3.4cm}|Y|}
\hline
\rowcolor{headerTabla}
\sffamily\textbf{Campo} & \sffamily\textbf{Valor} & \sffamily\textbf{Efecto} \\
\hline
\texttt{display} & \texttt{standalone} & Se abre sin barra de navegador, como app nativa. \\ \hline
\rowcolor{grisTecnico}
\texttt{orientation} & \texttt{portrait-primary} & Orientación vertical preferente en móvil. \\ \hline
\texttt{start\_url} · \texttt{scope} & \pathbreak{/} & Punto de arranque y ámbito de la instalación. \\ \hline
\rowcolor{grisTecnico}
\texttt{theme\_color} & \texttt{\#A855F7} & Color de la barra de estado (sincronizado con la paleta). \\ \hline
\texttt{background\_color} & \texttt{\#FCFCFC} & Color de la pantalla de arranque. \\ \hline
\rowcolor{grisTecnico}
\texttt{icons} & 192 / 512 px \texttt{maskable} & Iconos adaptativos para la pantalla de inicio. \\ \hline
\end{tabularx}
\label{tab:pwa-manifest}

\paragraph{Capacidades nativas en iOS y Android.} Además del manifiesto, el documento
\pathbreak{index.html} declara las \textit{meta} \comando{apple-mobile-web-app-capable},
\comando{apple-mobile-web-app-status-bar-style} (\texttt{black-translucent}) y
\comando{apple-touch-icon}, junto con un \textit{viewport} con \texttt{viewport-fit=cover}
que respeta las zonas seguras (\textit{notch}) de los dispositivos modernos.

\begin{AvisoNota}
La PWA prioriza la \textbf{instalabilidad} y la integración nativa (icono, color de tema,
modo \textit{standalone}, zonas seguras) sobre el funcionamiento sin conexión completo:
el sistema es intrínsecamente conectado (consume la API REST y recibe eventos Realtime),
por lo que el énfasis está en la experiencia tipo app más que en una caché \textit{offline}
exhaustiva.
\end{AvisoNota}
       % 5.3 Sitemap y enrutamiento + PWA
% =============================================================================
% 5.4 ESTRATEGIA DE INTERFAZ DE USUARIO Y COMPONENTES AVANZADOS
% =============================================================================
\section{Estrategia de Interfaz de Usuario y Wireframes}
\label{sec:ui}

La capa visual de EthosHub equilibra accesibilidad, consistencia y riqueza funcional.
Esta sección documenta la base de estilización y las integraciones avanzadas —mapas,
edición de código y gráficas— que diferencian la experiencia.

\subsection{Primitivas accesibles (Radix UI) y estilización atómica (TailwindCSS)}
\label{subsec:radix_tailwind}

La capa visual combina \textbf{primitivas accesibles} con \textbf{estilización
utilitaria}. Los componentes de \pathbreak{components/ui} encapsulan primitivas de Radix
(diálogos, menús, \textit{tooltips}, \textit{slots}) garantizando accesibilidad de origen
(gestión de foco, navegación por teclado y atributos ARIA), mientras que \textbf{TailwindCSS}
aplica estilos atómicos sin hojas de estilo ad-hoc. Las utilidades
\comando{class-variance-authority}, \comando{tailwind-merge} y \comando{clsx} gestionan
variantes de componentes de forma tipada y sin colisiones de clases.

\paragraph{Sistema de \textit{tokens} de diseño y tema.} Los colores, espaciados y radios
no se codifican de forma literal: se definen como \textbf{\textit{design tokens}}
(variables CSS) que TailwindCSS consume. El color primario de marca (\texttt{\#A855F7})
se declara una sola vez y se propaga al \texttt{theme\_color} de la PWA
(sección~\ref{subsec:pwa}, pág.~\pageref{subsec:pwa}). Este enfoque centralizado habilita
el cambio de tema (claro/oscuro) y la coherencia visual sin duplicar reglas de estilo.

\paragraph{Accesibilidad de origen.} Las primitivas de Radix aportan accesibilidad
\textbf{por construcción}: gestión de foco al abrir y cerrar diálogos, navegación completa
por teclado, atributos ARIA correctos y soporte de lectores de pantalla. La fuente
corporativa (\textbf{Sora}) y el respeto a las zonas seguras del dispositivo
(\texttt{viewport-fit=cover}) completan una experiencia inclusiva y consistente entre
plataformas.

\subsection{Componentes avanzados integrados}
\label{subsec:comp_avanzados}

La tabla~\ref{tab:libs-fe} cataloga las librerías especializadas integradas y el módulo
donde se consumen. A continuación se detallan las tres integraciones de mayor complejidad.

\vspace{0.2cm}
\noindent
\renewcommand{\arraystretch}{1.3}
\fontsize{8.5}{10.2}\selectfont\sffamily
\begin{tabularx}{\textwidth}{|L{3.0cm}|Y|L{3.2cm}|}
\hline
\rowcolor{headerTabla}
\sffamily\textbf{Librería} & \sffamily\textbf{Uso} & \sffamily\textbf{Módulo} \\
\hline
Monaco Editor & Edición de código LaTeX del CV con resaltado. & CV Studio \\
\hline
\rowcolor{grisTecnico}
Recharts & Gráficas: \textit{funnel}, radar de \textit{skills}, KPIs. & Dashboard reclutador \\
\hline
Leaflet · react-leaflet & Mapas interactivos y geolocalización. & Experiencia / preferencias \\
\hline
\rowcolor{grisTecnico}
@dnd-kit & \textit{Drag \& drop} del \textit{pipeline} Kanban. & Talent Discovery \\
\hline
Framer Motion & Animaciones de transición. & Transversal \\
\hline
\rowcolor{grisTecnico}
html2pdf.js & Exportación a PDF en cliente. & Portafolio / CV \\
\hline
react-webcam & Captura de foto desde la cámara. & Adjuntos de chat \\
\hline
\rowcolor{grisTecnico}
Sonner & Notificaciones \textit{toast}. & Transversal \\
\hline
\end{tabularx}
\label{tab:libs-fe}

\subsubsection{Navegación con mapas: Leaflet, OpenStreetMap y geocodificación}
\label{subsubsec:mapas}

La gestión de experiencia laboral (\pathbreak{/dashboard/experience}) integra un mapa
interactivo basado en \textbf{react-leaflet} sobre teselas de \textbf{OpenStreetMap}
servidas por CARTO. El usuario no escribe coordenadas: las \textbf{selecciona en el mapa}
o busca una dirección, y el sistema resuelve el resto mediante geocodificación. La
figura~\ref{fig:flow-mapa} formaliza este flujo bidireccional.

\vspace{0.3cm}
\begin{figure}[h!]
\centering
\begin{tikzpicture}[node distance=0.7cm and 1.5cm,
    box/.style={rectangle, rounded corners=2pt, draw, font=\sffamily\scriptsize, align=center,
        minimum width=2.8cm, minimum height=0.85cm, inner sep=3pt},
    rel/.style={-{Stealth[length=2mm]}, semithick}]
    \node[box, fill=c4Persona!25] (user) {Usuario};
    \node[box, fill=c4Componente!40, right=of user] (map) {\texttt{MapContainer}\\\scriptsize click / marcador};
    \node[box, fill=colorAcento!28, draw=colorAcento!80!black, right=of map] (nom) {Nominatim\\\scriptsize (OSM)};
    \node[box, fill=c4Componente!40, below=of map] (form) {Formulario\\\scriptsize lat · lng · dirección};

    \draw[rel] (user) -- node[c4etiqueta, above]{\scriptsize clic} (map);
    \draw[rel] (map) -- node[c4etiqueta, above]{\scriptsize reverse(lat,lng)} (nom);
    \draw[rel] (nom) to[bend right=20] node[c4etiqueta, right]{\scriptsize dirección} (form);
    \draw[rel] (map) -- node[c4etiqueta, left]{\scriptsize coordenadas} (form);
    \draw[rel] (user) to[bend right=30] node[c4etiqueta, below]{\scriptsize búsqueda de texto} (nom);
\end{tikzpicture}
\caption{Flujo de geolocalización bidireccional con Leaflet y Nominatim (OpenStreetMap).}
\label{fig:flow-mapa}
\end{figure}

El componente realiza dos operaciones de geocodificación contra el servicio
\textbf{Nominatim} de OpenStreetMap:

\begin{itemize}[noitemsep]
    \item \textbf{Geocodificación inversa} (\texttt{reverse}): al hacer clic en el mapa,
    convierte las coordenadas \texttt{(lat, lng)} en una dirección textual legible.
    \item \textbf{Geocodificación directa} (\texttt{search}): al escribir una dirección,
    obtiene hasta 6 candidatos con sus coordenadas para situar el marcador.
\end{itemize}

\begin{lstlisting}[language=JavaScript, caption={Geocodificación inversa con Nominatim (\texttt{pages/dashboard/ExperiencePage.tsx}).}, label={lst:nominatim}]
async function nominatimReverse(lat: number, lng: number): Promise<string> {
  const res = await fetch(
    `https://nominatim.openstreetmap.org/reverse?lat=${lat}&lon=${lng}&format=json`
  );
  const data = await res.json();
  return data.display_name ?? '';
}
\end{lstlisting}

\begin{NotaTecnica}
Los iconos de marcador de Leaflet se resuelven mediante \texttt{new URL(..., import.meta.url)}
para que Vite los empaquete correctamente en producción; sin esta resolución, el
empaquetado rompería las rutas relativas a las imágenes del marcador. Además, las
coordenadas almacenadas enlazan a \pathbreak{openstreetmap.org} para abrir la ubicación en
una pestaña externa.
\end{NotaTecnica}

\subsubsection{Edición de código: Monaco Editor en el CV Studio}
\label{subsubsec:monaco}

El CV Studio embebe \textbf{Monaco Editor} (el motor de edición de VS Code) para escribir
el código \textbf{LaTeX} del currículum con resaltado de sintaxis, plegado de bloques y
numeración de líneas. El contenido se sincroniza con el \texttt{cvStudioStore}
(sección~\ref{subsec:stores}, pág.~\pageref{subsec:stores}) y se envía al backend para su
compilación a PDF (capítulo~\ref{ch:backend}, figura~\ref{fig:seq-cv},
pág.~\pageref{fig:seq-cv}).

El CV Studio es el componente cliente con mayor densidad funcional: combina el editor, un
panel de conversación con la IA y una vista previa del PDF compilado. El
\texttt{cvStudioService} expone las operaciones \comando{listDocuments},
\comando{createDocument}, \comando{updateDocument}, \comando{deleteDocument},
\comando{callAiAssist} (asistencia generativa, que conserva el historial de la
conversación) y \comando{compileLatexToPdf} (que recibe un \texttt{Blob} binario con el
PDF). La tabla~\ref{tab:cvstudio} relaciona cada interacción del usuario con su llamada al
backend.

\vspace{0.2cm}
\noindent
\renewcommand{\arraystretch}{1.3}
\fontsize{8.5}{10.2}\selectfont\sffamily
\begin{tabularx}{\textwidth}{|L{3.6cm}|L{3.6cm}|Y|}
\hline
\rowcolor{headerTabla}
\sffamily\textbf{Interacción} & \sffamily\textbf{Función del servicio} & \sffamily\textbf{Endpoint} \\
\hline
Listar mis CV & \texttt{listDocuments()} & \pathbreak{GET /cv-documents} \\ \hline
\rowcolor{grisTecnico}
Guardar cambios & \texttt{updateDocument(id)} & \pathbreak{PUT /cv-documents/\{id\}} \\ \hline
Pedir ayuda IA & \texttt{callAiAssist()} & \pathbreak{POST /cv-documents/ai-assist} \\ \hline
\rowcolor{grisTecnico}
Generar PDF & \texttt{compileLatexToPdf()} & \pathbreak{POST /cv-documents/compile-latex} \\ \hline
\end{tabularx}
\label{tab:cvstudio}

\begin{AvisoNota}
La compilación devuelve un \texttt{Blob} de tipo \texttt{application/pdf}, que el cliente
convierte en una URL de objeto (\texttt{URL.createObjectURL}) para mostrar la vista previa
sin recargar la página. Si el código LaTeX es inválido, el backend responde \texttt{422}
(capítulo~\ref{ch:backend}, tabla~\ref{tab:http-status}, pág.~\pageref{tab:http-status}) y
el cliente muestra el error sin perder el contenido del editor.
\end{AvisoNota}

\subsubsection{Visualización de datos: Recharts en el dashboard del reclutador}
\label{subsubsec:recharts}

El \textit{dashboard} del reclutador usa \textbf{Recharts} para tres visualizaciones: un
\textbf{embudo} (\textit{funnel}) del \textit{pipeline} de candidatos, un \textbf{radar}
de cobertura de \textit{skills} por candidato y \textbf{KPIs} de actividad. Estas gráficas
consumen los datos agregados expuestos por \texttt{dashboardService}
(sección~\ref{subsec:apiclient}, pág.~\pageref{subsec:apiclient}).

\subsubsection{Tablero Kanban de talento con \textit{drag \& drop} (@dnd-kit)}
\label{subsubsec:kanban}

El descubrimiento de talento (\pathbreak{/recruiter/talent-discovery}) organiza a los
candidatos en un \textbf{tablero Kanban} con columnas que representan fases del
\textit{pipeline} de selección. El reclutador arrastra tarjetas entre columnas mediante
\textbf{@dnd-kit}, una librería accesible de \textit{drag \& drop} que respeta el teclado
y los lectores de pantalla. Cada movimiento persiste el cambio de fase a través de los
\textit{endpoints} de \textit{likes}/fases del backend, que aplican las políticas RLS
asimétricas descritas en el capítulo~\ref{ch:basedatos}
(sección~\ref{subsec:politicas_rls}, pág.~\pageref{subsec:politicas_rls}).

\subsection{Wireframe de referencia: espacio de trabajo del reclutador}
\label{subsec:wireframe}

La figura~\ref{fig:wireframe} presenta un \textit{wireframe} de baja fidelidad del
\textit{dashboard} del reclutador, que integra las visualizaciones Recharts, el tablero
Kanban y el asistente de IA en una sola vista.

\vspace{0.3cm}
\begin{figure}[h!]
\centering
\begin{tikzpicture}[
    pane/.style={rectangle, draw=grisISO!70, thick, fill=grisTecnico, font=\sffamily\scriptsize,
        align=center, inner sep=4pt},
    lbl/.style={font=\sffamily\scriptsize\bfseries, text=colorPrimario}]
    % Marco general
    \draw[draw=grisISO, very thick, rounded corners=3pt] (-0.2,-0.2) rectangle (11.2,6.4);
    % Barra superior
    \node[pane, minimum width=11cm, minimum height=0.7cm] at (5.5,5.9) {\lbl Top bar · logo · selector de idioma · perfil};
    % Sidebar (móvil colapsable)
    \node[pane, minimum width=2.0cm, minimum height=4.6cm] at (1.0,3.0) {Nav\\(colapsable\\en móvil)};
    % KPIs
    \node[pane, minimum width=8.0cm, minimum height=0.9cm] at (6.6,5.0) {\lbl KPIs · candidatos · vistas · mensajes};
    % Funnel + Radar
    \node[pane, minimum width=3.9cm, minimum height=1.9cm] at (4.6,3.5) {Embudo\\(Recharts)};
    \node[pane, minimum width=3.9cm, minimum height=1.9cm] at (8.7,3.5) {Radar de skills\\(Recharts)};
    % Kanban
    \node[pane, minimum width=5.9cm, minimum height=1.6cm] at (5.6,1.5) {Kanban (@dnd-kit)\\\scriptsize Nuevo · Revisando · Contactado};
    % AI insights
    \node[pane, minimum width=1.9cm, minimum height=1.6cm, fill=colorAcento!18] at (9.7,1.5) {Insights\\IA};
\end{tikzpicture}
\caption{Wireframe de baja fidelidad del \textit{dashboard} del reclutador.}
\label{fig:wireframe}
\end{figure}
            % 5.4 Estrategia de interfaz y wireframes
% =============================================================================
% 5.5 GESTIÓN DE ESTADO GLOBAL Y LOCAL (ZUSTAND)
% =============================================================================
\section{Gestión de Estado Global y Local (Zustand)}
\label{sec:zustand}

El estado de la aplicación se gestiona con \textbf{Zustand}, una librería minimalista que
evita el \textit{boilerplate} de soluciones más pesadas. El diseño reparte el estado en
\textit{stores} independientes por dominio, en lugar de un único almacén monolítico.

\subsection{Arquitectura de \textit{stores} por dominios}
\label{subsec:stores}

Cada \textit{store} encapsula el estado y las acciones de un dominio concreto, lo que
limita el alcance de las re-renderizaciones y mantiene el código cohesionado. La
tabla~\ref{tab:stores} enumera los \textit{stores} reales del directorio \pathbreak{store/}.

\vspace{0.2cm}
\noindent
\renewcommand{\arraystretch}{1.3}
\fontsize{8.5}{10.2}\selectfont\sffamily
\begin{tabularx}{\textwidth}{|L{4.0cm}|Y|}
\hline
\rowcolor{headerTabla}
\sffamily\textbf{Store} & \sffamily\textbf{Responsabilidad} \\
\hline
\texttt{authStore} & Sesión, token, rol e identidad; persistencia e hidratación de idioma. \\
\hline
\rowcolor{grisTecnico}
\texttt{notificationsStore} & Notificaciones reactivas del usuario (con eventos Realtime). \\
\hline
\texttt{projectsStore} · \texttt{skillsStore} & Estado de proyectos y habilidades del perfil. \\
\hline
\rowcolor{grisTecnico}
\texttt{cvStudioStore} & Estado del editor de CV y la conversación con la IA. \\
\hline
\texttt{portfolioStore} · \texttt{connectionsStore} & Ajustes del portafolio y conexiones externas. \\
\hline
\rowcolor{grisTecnico}
\texttt{preferencesStore} · \texttt{uiStore} & Preferencias del usuario y estado efímero de la interfaz. \\
\hline
\texttt{resetAllStores} & Utilidad que reinicia todos los \textit{stores} al cerrar sesión. \\
\hline
\end{tabularx}
\label{tab:stores}

\subsection{Mecanismos de persistencia en el cliente}
\label{subsec:persistencia}

La gestión de sesión es un punto de diseño cuidadoso. El \comando{authStore} emplea el
\textit{middleware} \comando{persist} de Zustand y distingue dos almacenamientos del
navegador con propósitos distintos, descritos en la tabla~\ref{tab:persistencia}.

\clearpage
\vspace{0.2cm}
\noindent
\renewcommand{\arraystretch}{1.3}
\fontsize{8.5}{10.2}\selectfont\sffamily
\begin{tabularx}{\textwidth}{|L{3.2cm}|L{3.8cm}|Y|}
\hline
\rowcolor{headerTabla}
\sffamily\textbf{Almacén} & \sffamily\textbf{Clave} & \sffamily\textbf{Contenido y motivo} \\
\hline
\texttt{sessionStorage} & \texttt{ethoshub\_access\_token} & Token y expiración; \textbf{por pestaña}, habilita sesiones múltiples aisladas. \\ \hline
\rowcolor{grisTecnico}
\texttt{localStorage} & \texttt{ethoshub\_auth} & Porción no sensible del \textit{store}, para hidratar la identidad entre recargas. \\ \hline
\end{tabularx}
\label{tab:persistencia}

\paragraph{Ciclo de vida de la sesión.} La sesión atraviesa un ciclo de vida bien
definido —desde la hidratación al arrancar la app hasta la limpieza al cerrar sesión—,
representado como máquina de estados en la figura~\ref{fig:fsm-sesion}.

\vspace{0.3cm}
\begin{figure}[h!]
\centering
\begin{tikzpicture}[node distance=1.1cm and 1.9cm,
    estado/.style={rectangle, rounded corners=8pt, draw=c4Sistema!70!black, fill=c4Componente!30,
        font=\sffamily\scriptsize\bfseries, align=center, minimum width=2.1cm, minimum height=0.75cm},
    ini/.style={circle, fill=flujoInicio, minimum size=0.4cm, inner sep=0pt},
    rel/.style={-{Stealth[length=2mm]}, semithick}]
    \node[ini] (s) {};
    \node[estado, right=1.0cm of s] (hidr) {Hidratando\\\scriptsize (localStorage)};
    \node[estado, right=of hidr] (anon) {Anónimo};
    \node[estado, right=of anon] (auth) {Autenticado\\\scriptsize (token activo)};
    \node[estado, below=of auth] (exp) {Expirado};

    \draw[rel] (s) -- (hidr);
    \draw[rel] (hidr) -- node[c4etiqueta, above]{sin token} (anon);
    \draw[rel] (anon) -- node[c4etiqueta, above]{login} (auth);
    \draw[rel] (auth) -- node[c4etiqueta, right]{exp vencido} (exp);
    \draw[rel] (exp) to[bend left=30] node[c4etiqueta, below]{auth:unauthorized} (anon);
    \draw[rel] (auth) to[bend right=45] node[c4etiqueta, above]{logout} (anon);
\end{tikzpicture}
\caption{Máquina de estados del ciclo de vida de la sesión en el cliente.}
\label{fig:fsm-sesion}
\end{figure}

\paragraph{Cierre de sesión coordinado entre pestañas.} El \textit{logout} es
\textbf{global}: \comando{logout()} limpia el \texttt{sessionStorage} y el
\texttt{localStorage} y emite un evento por un \comando{BroadcastChannel('ethoshub\_auth')},
de modo que todas las pestañas abiertas cierran sesión de forma sincronizada. La
figura~\ref{fig:flow-logout} ilustra la propagación.

\vspace{0.3cm}
\begin{figure}[h!]
\centering
\begin{tikzpicture}[node distance=0.7cm and 1.6cm,
    box/.style={rectangle, rounded corners=2pt, draw, font=\sffamily\scriptsize, align=center,
        minimum width=2.6cm, minimum height=0.8cm, inner sep=3pt},
    rel/.style={-{Stealth[length=2mm]}, semithick}]
    \node[box, fill=bgFail!85, text=white, draw=red!70!black] (lo) {Pestaña A:\\\texttt{logout()}};
    \node[box, fill=c4Componente!40, right=of lo] (clear) {limpia storage\\\scriptsize session + local};
    \node[box, fill=colorAcento!28, draw=colorAcento!80!black, right=of clear] (bc) {BroadcastChannel\\\scriptsize ethoshub\_auth};
    \node[box, fill=c4Componente!40, below=of bc] (tabB) {Pestaña B:\\\scriptsize cierra sesión};
    \node[box, fill=c4Componente!40, below=of clear] (tabC) {Pestaña C:\\\scriptsize cierra sesión};

    \draw[rel] (lo) -- (clear);
    \draw[rel] (clear) -- (bc);
    \draw[rel] (bc) -- node[c4etiqueta, right]{\scriptsize evento} (tabB);
    \draw[rel] (bc) to[bend left=10] node[c4etiqueta, above]{\scriptsize evento} (tabC);
\end{tikzpicture}
\caption{Cierre de sesión global propagado entre pestañas vía \texttt{BroadcastChannel}.}
\label{fig:flow-logout}
\end{figure}

\begin{NotaTecnica}
La elección de \texttt{sessionStorage} para el token —y no \texttt{localStorage}— es
deliberada: al estar aislado por pestaña, un usuario puede mantener \textbf{dos sesiones
simultáneas} (p.~ej., profesional y reclutador) en pestañas distintas sin que una
sobrescriba a la otra. La porción persistida en \texttt{localStorage} es explícitamente
\textbf{no sensible} (no contiene el token), solo datos de identidad para hidratación.
\end{NotaTecnica}

\subsection{Estado reactivo en tiempo real: \texttt{notificationsStore}}
\label{subsec:rt_cliente}

El \comando{notificationsStore} es \textbf{reactivo}: no consulta el servidor por
\textit{polling}, sino que se actualiza ante eventos \textit{push} de \textbf{Supabase
Realtime}. Cuando el backend persiste una notificación (capítulo~\ref{ch:backend},
sección~\ref{sec:realtime}, pág.~\pageref{sec:realtime}), Supabase emite el cambio al
cliente suscrito y el \textit{store} incrementa el contador y reordena la lista sin
recargar la página. La figura~\ref{fig:flow-rt-cliente} sintetiza la suscripción.

\vspace{0.3cm}
\begin{figure}[h!]
\centering
\begin{tikzpicture}[node distance=0.7cm and 1.6cm,
    box/.style={rectangle, rounded corners=2pt, draw, font=\sffamily\scriptsize, align=center,
        minimum width=2.7cm, minimum height=0.8cm, inner sep=3pt},
    rel/.style={-{Stealth[length=2mm]}, semithick}]
    \node[box, fill=colorAcento!28, draw=colorAcento!80!black] (rt) {Supabase\\Realtime};
    \node[box, fill=c4Componente!45, right=of rt] (store) {notifications\\Store};
    \node[box, fill=c4Componente!40, right=of store] (ui) {Campana\\\scriptsize + badge contador};

    \draw[rel] (rt) -- node[c4etiqueta, above]{\scriptsize evento INSERT} (store);
    \draw[rel] (store) -- node[c4etiqueta, above]{\scriptsize re-render} (ui);
\end{tikzpicture}
\caption{Consumo reactivo de notificaciones en tiempo real en el cliente.}
\label{fig:flow-rt-cliente}
\end{figure}

El usuario marca las notificaciones como leídas individualmente o en bloque; estas
acciones llaman a \\ \texttt{notificationsService} (\pathbreak{/api/v1/notifications}), y el
contador del \textit{badge} se recalcula a partir del estado del \textit{store}.
       % 5.5 Gestión de estado (Zustand)
% =============================================================================
% 5.6 CICLO DE VIDA Y CONSUMO ASÍNCRONO DE SERVICIOS
% =============================================================================
\section{Ciclo de Vida y Consumo Asíncrono de Servicios}
\label{sec:servicios}

Todo el consumo de la API REST del backend (capítulo~\ref{ch:backend},
pág.~\pageref{ch:backend}) se canaliza a través de una única instancia de Axios
preconfigurada. Esta sección documenta su configuración, los interceptores que inyectan el
JWT y gestionan las sesiones expiradas, y el catálogo de servicios de dominio.

\subsection{Configuración del cliente base (\texttt{apiClient.ts})}
\label{subsec:apiclient}

El cliente HTTP se centraliza en \comando{apiClient.ts}, que todos los servicios de
\pathbreak{shared/services} reutilizan. Sus parámetros clave figuran en la
tabla~\ref{tab:apiclient}.

\vspace{0.2cm}
\noindent
\renewcommand{\arraystretch}{1.3}
\fontsize{8.5}{10.2}\selectfont\sffamily
\begin{tabularx}{\textwidth}{|L{3.4cm}|L{3.2cm}|Y|}
\hline
\rowcolor{headerTabla}
\sffamily\textbf{Parámetro} & \sffamily\textbf{Valor} & \sffamily\textbf{Motivo} \\
\hline
\texttt{baseURL} & \pathbreak{/api} o \texttt{VITE\_API\_URL} & Prefijo común; configurable por entorno. \\ \hline
\rowcolor{grisTecnico}
\texttt{timeout} & \texttt{8000} ms & Aborta peticiones colgadas y libera el indicador de carga. \\ \hline
\texttt{headers} & \texttt{Content-Type: application/json} & Contrato JSON por defecto. \\ \hline
\rowcolor{grisTecnico}
\texttt{withCredentials} & \texttt{true} & Envía \textit{cookies} de credenciales (token vía \textit{cookie}). \\ \hline
\end{tabularx}
\label{tab:apiclient}

\begin{NotaTecnica}
En desarrollo, Vite expone un \textit{proxy} que reenvía \pathbreak{/api}, \pathbreak{/auth}
y \pathbreak{/v1} a \pathbreak{http://localhost:8080} (el backend Spring), evitando
problemas de CORS durante el desarrollo local. En producción no hay \textit{proxy}: el
mismo \textit{JAR} sirve la SPA y la API bajo el mismo origen.
\end{NotaTecnica}

La figura~\ref{fig:pipeline-cliente} muestra el recorrido completo de una petición desde
una vista hasta el backend y de vuelta, atravesando los interceptores que se documentan en
los apartados siguientes. Es el espejo, en el cliente, del \textit{pipeline} de servidor
de la figura~\ref{fig:pipeline-http} (pág.~\pageref{fig:pipeline-http}).

\vspace{0.3cm}
\begin{figure}[h!]
\centering
\begin{tikzpicture}[node distance=0.5cm and 0.85cm,
    etapa/.style={rectangle, rounded corners=3pt, draw=c4Sistema!70!black, fill=c4Componente!30,
        font=\sffamily\scriptsize, align=center, minimum width=2.2cm, minimum height=0.8cm, inner sep=2pt},
    intc/.style={etapa, fill=colorAcento!25, draw=colorAcento!80!black},
    rel/.style={-{Stealth[length=2mm]}, semithick, draw=grisISO}]
    \node[etapa, fill=c4Persona!25] (view) {Vista / hook};
    \node[etapa, right=of view] (svc) {Servicio\\\scriptsize de dominio};
    \node[intc, right=of svc] (reqi) {Interceptor\\\scriptsize petición (JWT)};
    \node[etapa, fill=c4Datos!22, right=of reqi] (be) {Backend\\\scriptsize REST};
    \node[intc, below=0.8cm of be] (respi) {Interceptor\\\scriptsize respuesta};
    \node[etapa, fill=c4Persona!25, left=of respi] (store) {Store / UI};

    \draw[rel] (view) -- (svc);
    \draw[rel] (svc) -- (reqi);
    \draw[rel] (reqi) -- node[c4etiqueta, above]{\scriptsize Bearer} (be);
    \draw[rel] (be) -- (respi);
    \draw[rel] (respi) -- node[c4etiqueta, above]{\scriptsize 401/403?} (store);
\end{tikzpicture}
\caption{Pipeline de una petición en el cliente, con sus interceptores.}
\label{fig:pipeline-cliente}
\end{figure}

\subsection{Interceptor de peticiones: inyección del JWT}
\label{subsec:interceptor_req}

El interceptor de petición lee el token (de \texttt{sessionStorage}, con \textit{fallback}
a \texttt{localStorage}) y lo inyecta dinámicamente como cabecera \comando{Authorization:
Bearer <token>}, de modo que ningún servicio gestione el token manualmente.

\begin{lstlisting}[language=JavaScript, caption={Interceptor de petición: inyección del JWT (\texttt{apiClient.ts}).}, label={lst:req-interceptor}]
function readToken(): string | null {
  return sessionStorage.getItem(ACCESS_TOKEN_KEY)
      ?? localStorage.getItem(ACCESS_TOKEN_KEY);
}

apiClient.interceptors.request.use((config) => {
  const token = readToken();
  if (token) config.headers.Authorization = `Bearer ${token}`;
  return config;
});
\end{lstlisting}

\subsection{Interceptor de respuestas: manejo de sesiones expiradas}
\label{subsec:interceptor_resp}

El interceptor de respuesta concentra la \textbf{lógica de detección de sesión expirada}.
Ante un \texttt{401} o \texttt{403}, no asume ciegamente que la sesión caducó:
\textbf{decodifica el \textit{claim} \texttt{exp} del JWT} y lo contrasta con la marca de
expiración almacenada. Solo si el token está efectivamente expirado emite el evento
\comando{auth:unauthorized}, que dispara el cierre de sesión y la redirección. La
figura~\ref{fig:flow-401} formaliza esta decisión.

\vspace{0.3cm}
\begin{figure}[h!]
\centering
\begin{tikzpicture}[node distance=0.75cm and 1.4cm]
    \node[flujoInicio] (resp) {Respuesta HTTP};
    \node[flujoDecision, below=of resp] (d1) {¿401 / 403?};
    \node[flujoProceso, right=2.2cm of d1] (ok) {Propagar al servicio};
    \node[flujoDecision, below=1.0cm of d1] (d2) {¿exp vencido\\o token ausente?};
    \node[flujoProceso, right=2.2cm of d2] (keep) {Mantener sesión\\(error puntual)};
    \node[flujoInicio, fill=bgFail, draw=red!70!black, below=1.0cm of d2] (evt) {Emitir \texttt{auth:unauthorized}};

    \draw[flujoFlecha] (resp) -- (d1);
    \draw[flujoFlecha] (d1) -- node[c4etiqueta, above]{no} (ok);
    \draw[flujoFlecha] (d1) -- node[c4etiqueta, left]{sí} (d2);
    \draw[flujoFlecha] (d2) -- node[c4etiqueta, above]{no} (keep);
    \draw[flujoFlecha] (d2) -- node[c4etiqueta, left]{sí} (evt);
\end{tikzpicture}
\caption{Decisión del interceptor de respuesta ante errores 401/403.}
\label{fig:flow-401}
\end{figure}

\clearpage
\begin{lstlisting}[language=JavaScript, caption={Interceptor de respuesta: distinción entre expiración real y error puntual (\texttt{apiClient.ts}).}, label={lst:resp-interceptor}]
apiClient.interceptors.response.use(
  (response) => response,
  (error: AxiosError) => {
    const status = error.response?.status;
    if (status === 401 || status === 403) {
      const token     = readToken();
      const expiresAt = sessionStorage.getItem(EXPIRES_AT_KEY)
                     ?? localStorage.getItem(EXPIRES_AT_KEY);
      const isExpiredByStore = !token
        || (!!expiresAt && Date.now() > Number(expiresAt) * 1000);
      const isExpiredByJwt   = token ? isJwtExpired(token) : true;
      if (isExpiredByStore || isExpiredByJwt) {
        window.dispatchEvent(new CustomEvent('auth:unauthorized'));
      }
    }
    return Promise.reject(error);
  },
);
\end{lstlisting}

\begin{AvisoNota}
Esta distinción evita cerrar la sesión del usuario por un \texttt{403} legítimo (p.~ej.,
acceder a un recurso ajeno protegido por las políticas RLS del capítulo~\ref{ch:basedatos},
pág.~\pageref{ch:basedatos}) o por un \texttt{401} transitorio: solo se fuerza el
\textit{logout} cuando el JWT está \textbf{realmente} caducado, mejorando la experiencia y
evitando expulsiones espurias. El interceptor también libera el indicador de carga global
ante \textit{timeouts} (\texttt{ECONNABORTED}) y errores de servidor ($\ge$500).
\end{AvisoNota}

\subsection{Catálogo de servicios de dominio}
\label{subsec:servicios_catalogo}

La tabla~\ref{tab:servicios-fe} relaciona los servicios de \pathbreak{shared/services} con
su responsabilidad; todos construidos sobre \texttt{apiClient} y devolviendo \textit{tipos}
de TypeScript derivados del contrato \texttt{ApiResponse<T>} del backend.

\vspace{0.2cm}
\noindent
\renewcommand{\arraystretch}{1.3}
\fontsize{8.5}{10.2}\selectfont\sffamily
\begin{tabularx}{\textwidth}{|L{4.4cm}|Y|}
\hline
\rowcolor{headerTabla}
\sffamily\textbf{Servicio} & \sffamily\textbf{Responsabilidad} \\
\hline
\texttt{authService} & Registro, \textit{login}, sincronización OAuth, refresco de sesión. \\
\hline
\rowcolor{grisTecnico}
\texttt{cvStudioService} & Operaciones del CV: edición, asistencia IA, compilación a PDF. \\
\hline
\texttt{dashboardService} & Métricas y analíticas del reclutador. \\
\hline
\rowcolor{grisTecnico}
\texttt{educationService} · \texttt{experienceService} & Formación y experiencia laboral (con geolocalización). \\
\hline
\texttt{fileService} & Subida de archivos e integración con Drive. \\
\hline
\rowcolor{grisTecnico}
\texttt{notificationsService} · \texttt{preferencesService} & Notificaciones y preferencias del usuario. \\
\hline
\texttt{portfolioService} · \texttt{companyProfileService} & Portafolio público y perfil de empresa. \\
\hline
\rowcolor{grisTecnico}
\texttt{adminService} & Operaciones del panel de administración. \\
\hline
\end{tabularx}
\label{tab:servicios-fe}

\clearpage
\subsection{Resiliencia del cliente y manejo de errores}
\label{subsec:resiliencia}

El cliente aplica tres mecanismos de resiliencia que, combinados, evitan que un error
puntual degrade toda la experiencia. La tabla~\ref{tab:resiliencia} los resume.

\vspace{0.2cm}
\noindent
\renewcommand{\arraystretch}{1.3}
\fontsize{8.5}{10.2}\selectfont\sffamily
\begin{tabularx}{\textwidth}{|L{3.6cm}|Y|}
\hline
\rowcolor{headerTabla}
\sffamily\textbf{Mecanismo} & \sffamily\textbf{Función} \\
\hline
\texttt{RouteErrorBoundary} & Captura errores de renderizado por ruta y muestra una vista de error en lugar de una pantalla en blanco. \\ \hline
\rowcolor{grisTecnico}
\texttt{Suspense} + \textit{skeleton} & Mantiene la interfaz responsiva mientras carga el \textit{chunk} diferido de una página. \\ \hline
\texttt{AbortController} & Cancela peticiones en vuelo al desmontar la vista, evitando condiciones de carrera. \\ \hline
\rowcolor{grisTecnico}
\textit{timeout} 8\,s + \textit{toast} & Aborta peticiones colgadas y notifica al usuario con un \textit{toast} (Sonner). \\ \hline
\end{tabularx}
\label{tab:resiliencia}

\paragraph{Cancelación de peticiones.} El módulo exporta \comando{createAbortController()},
un \textit{factory} de \texttt{AbortController} que las vistas usan para \textbf{cancelar}
peticiones en vuelo al desmontarse (p.~ej., al navegar fuera de una página de búsqueda),
evitando actualizaciones de estado sobre componentes ya desmontados y condiciones de
carrera. El interceptor de respuesta distingue además las cancelaciones
(\texttt{isCancel}) de los errores reales, de modo que abortar una petición no se
contabiliza como fallo.

\paragraph{Separación de responsabilidades en el consumo.} El diseño del consumo de
servicios respeta una separación de responsabilidades nítida que conviene subrayar como
cierre. La \textbf{vista} (o su \textit{custom hook}) decide \textit{cuándo} pedir datos y
\textit{qué} hacer con ellos, pero no conoce el detalle del transporte. El \textbf{servicio
de dominio} traduce esa intención a una llamada concreta y devuelve tipos de TypeScript, sin
saber nada de tokens ni de sesiones. El \textbf{interceptor} inyecta el JWT y gobierna las
sesiones expiradas de forma transversal, sin que ningún servicio lo gestione. Y el
\textbf{store} mantiene el estado resultante, notificando a las vistas suscritas. Esta
estratificación —vista, servicio, interceptor, store— es la contraparte cliente de las
capas Controller~/~Service~/~Repository del backend (capítulo~\ref{ch:backend},
sección~\ref{sec:clases}, pág.~\pageref{sec:clases}): en ambos extremos, cada capa tiene una
única razón de cambio, lo que hace el sistema predecible y mantenible a lo largo de su
ciclo de vida, tal como persigue la norma ISO/IEC/IEEE~15289:2024.
     % 5.6 Consumo asíncrono de servicios
% =============================================================================
% 5.7 INTERNACIONALIZACIÓN, AUTENTICACIÓN OAUTH Y OPTIMIZACIÓN
% =============================================================================
\section{Internacionalización, Autenticación Federada y Optimización}
\label{sec:i18n_build}

Esta sección documenta tres aspectos transversales del frontend que completan el diseño
detallado: la internacionalización multilingüe, el flujo de autenticación federada OAuth
en el cliente, y las estrategias de optimización del rendimiento de carga.

\subsection{Internacionalización con i18next (es / en / pt)}
\label{subsec:i18n}

La aplicación es \textbf{trilingüe} —español, inglés y portugués— mediante
\textbf{i18next} y \textbf{react-i18next}. Las traducciones se organizan en \textbf{11
espacios de nombres} (\textit{namespaces}) por idioma, lo que segmenta los textos por
área funcional y evita un único fichero monolítico. La tabla~\ref{tab:i18n-ns} enumera los
\textit{namespaces}.

\clearpage
\vspace{0.2cm}
\noindent
\renewcommand{\arraystretch}{1.3}
\fontsize{8.5}{10.2}\selectfont\sffamily
\begin{tabularx}{\textwidth}{|L{3.0cm}|Y|}
\hline
\rowcolor{headerTabla}
\sffamily\textbf{Namespace} & \sffamily\textbf{Contenido} \\
\hline
\texttt{common} · \texttt{nav} & Textos transversales y de navegación. \\ \hline
\rowcolor{grisTecnico}
\texttt{auth} · \texttt{settings} & Pantallas de acceso y de preferencias. \\ \hline
\texttt{dashboard} · \texttt{recruiter} · \texttt{admin} & Espacios de trabajo por rol. \\ \hline
\rowcolor{grisTecnico}
\texttt{public} · \texttt{pages} & Vistas públicas y páginas estáticas. \\ \hline
\texttt{domain} · \texttt{components} & Terminología de dominio y componentes reutilizables. \\ \hline
\end{tabularx}
\label{tab:i18n-ns}

El idioma se resuelve con \comando{i18next-browser-languagedetector}, que infiere la
preferencia del navegador, y se sincroniza con el \comando{authStore} para persistir la
elección del usuario entre sesiones (hidratación de idioma, sección~\ref{subsec:stores},
pág.~\pageref{subsec:stores}). El backend respeta esta preferencia para localizar correos
y mensajes vía \pathbreak{/api/v1/preferences}.

Los componentes consumen las traducciones con el \textit{hook} \comando{useTranslation},
referenciando el \textit{namespace} y la clave, como ilustra el
listado~\ref{lst:i18n-uso}. La inicialización (listado~\ref{lst:i18n-init}) carga los 33
módulos de traducción (11 \textit{namespaces} $\times$ 3 idiomas) y registra el detector.

\begin{lstlisting}[language=JavaScript, caption={Inicialización de i18next con detector de idioma (\texttt{i18n/index.ts}).}, label={lst:i18n-init}]
i18n
  .use(LanguageDetector)
  .use(initReactI18next)
  .init({
    resources: { es, en, pt },     // 11 namespaces por idioma
    fallbackLng: 'es',
    interpolation: { escapeValue: false },
  });
\end{lstlisting}

\begin{lstlisting}[language=JavaScript, caption={Consumo de traducciones en un componente.}, label={lst:i18n-uso}]
const { t } = useTranslation('dashboard');
return <h1>{t('portfolio.title')}</h1>;  // -> "Mi portafolio" / "My portfolio"
\end{lstlisting}

\subsection{Flujo de autenticación federada (OAuth2) en el cliente}
\label{subsec:oauth}

Además del acceso con credenciales, EthosHub admite \textbf{OAuth2} (proveedores sociales
vía Supabase Auth). El cliente no maneja secretos de OAuth: delega el intercambio en
Supabase y solo recibe el resultado en una ruta de \textit{callback}. La
figura~\ref{fig:flow-oauth} ilustra el flujo.

\vspace{0.3cm}
\begin{figure}[h!]
\centering
\begin{tikzpicture}[node distance=0.7cm and 1.3cm,
    box/.style={rectangle, rounded corners=2pt, draw, font=\sffamily\scriptsize, align=center,
        minimum width=2.6cm, minimum height=0.85cm, inner sep=3pt},
    rel/.style={-{Stealth[length=2mm]}, semithick}]
    \node[box, fill=c4Persona!25] (user) {Usuario};
    \node[box, fill=c4Componente!40, right=of user] (login) {LoginPage\\\scriptsize ``Entrar con...''};
    \node[box, fill=colorAcento!28, draw=colorAcento!80!black, right=of login] (sb) {Supabase\\Auth (OAuth)};
    \node[box, fill=c4Componente!40, below=of sb] (cb) {OAuth2Callback\\Page};
    \node[box, fill=c4Componente!45, below=of login] (sync) {authService\\\scriptsize POST /oauth/sync};
    \node[box, fill=c4Sistema, text=white, left=of sync] (home) {Dashboard\\por rol};

    \draw[rel] (user) -- (login);
    \draw[rel] (login) -- node[c4etiqueta, above]{\scriptsize redirect} (sb);
    \draw[rel] (sb) -- node[c4etiqueta, right]{\scriptsize token} (cb);
    \draw[rel] (cb) -- (sync);
    \draw[rel] (sync) -- node[c4etiqueta, above]{\scriptsize sesión} (home);
\end{tikzpicture}
\caption{Flujo de autenticación federada OAuth2 en el cliente.}
\label{fig:flow-oauth}
\end{figure}

Tras el retorno, \comando{OAuth2CallbackPage} extrae el token de Supabase y llama a
\comando{authService} para sincronizar el perfil en el backend
(\pathbreak{POST /api/auth/oauth/sync}, capítulo~\ref{ch:backend}, tabla~\ref{tab:endpoints},
pág.~\pageref{tab:endpoints}), creando el perfil local si es la primera vez. El servidor
de desarrollo de Vite establece la cabecera \\
\comando{Cross-Origin-Opener-Policy: same-origin-allow-popups} para permitir el flujo de
ventana emergente de OAuth.

\subsection{Optimización del rendimiento de carga}
\label{subsec:optimizacion}

El rendimiento percibido se optimiza con tres técnicas combinadas, descritas en la
tabla~\ref{tab:optimizacion}.

\vspace{0.2cm}
\noindent
\renewcommand{\arraystretch}{1.3}
\fontsize{8.5}{10.2}\selectfont\sffamily
\begin{tabularx}{\textwidth}{|L{3.6cm}|Y|}
\hline
\rowcolor{headerTabla}
\sffamily\textbf{Técnica} & \sffamily\textbf{Efecto} \\
\hline
\textit{Code splitting} por ruta & Cada página se carga bajo demanda (\texttt{React.lazy} + \texttt{Suspense}). \\ \hline
\rowcolor{grisTecnico}
\textit{Manual chunks} (Vite) & Se aísla el \textit{vendor} (\texttt{react}, \texttt{react-dom}, \texttt{react-router-dom}) en un \textit{chunk} cacheable. \\ \hline
\textit{Optimize deps} & Se pre-empaquetan \texttt{leaflet} y \texttt{react-leaflet} para acelerar el arranque en desarrollo. \\ \hline
\rowcolor{grisTecnico}
\textit{Skeletons} de carga & \texttt{Suspense} muestra esqueletos en vez de pantallas en blanco durante la carga diferida. \\ \hline
\end{tabularx}
\label{tab:optimizacion}

\begin{NotaTecnica}
La división del \textit{vendor} en su propio \textit{chunk} es una optimización de caché:
las dependencias base (React y el enrutador) cambian con poca frecuencia, por lo que el
navegador puede reutilizar ese \textit{chunk} entre despliegues mientras solo se
re-descarga el código de aplicación modificado.
\end{NotaTecnica}    % 5.7 i18n, OAuth y optimización
% =============================================================================
% 5.8 CALIDAD DEL CÓDIGO Y HERRAMIENTAS DEL FRONTEND
% =============================================================================
\section{Calidad del Código y Herramientas del Frontend}
\label{sec:calidad_fe}

El diseño del frontend incorpora \textbf{garantías de calidad en tiempo de construcción}
que detectan errores antes del despliegue. Esta sección documenta el \textit{tooling} que
las hace cumplir; la estrategia integral de verificación y validación se aborda en el
capítulo~8.

\subsection{Seguridad de tipos con TypeScript estricto}
\label{subsec:typescript}

Todo el código está escrito en \textbf{TypeScript con modo estricto}. El \textit{script}
de construcción ejecuta primero \comando{tsc} (verificación de tipos) y solo si pasa
procede con \comando{vite build}; un error de tipos \textbf{aborta el empaquetado}. Los
tipos de respuesta de los servicios derivan del contrato \texttt{ApiResponse<T>} del
backend (capítulo~\ref{ch:backend}, listado~\ref{lst:apiresponse},
pág.~\pageref{lst:apiresponse}), de modo que un cambio incompatible en la forma de los
datos del servidor se detecta como error de compilación en el cliente.

\subsection{Linting y convenciones (ESLint)}
\label{subsec:eslint}

El proyecto aplica \textbf{ESLint 10} con los \textit{plugins} de React Hooks y React
Refresh, que verifican las reglas de los \textit{hooks} (dependencias completas, llamadas
en el nivel superior) y la compatibilidad con la recarga en caliente. El \textit{script}
\comando{npm run lint} reporta las directivas no usadas y limita las advertencias.

\subsection{Pruebas con Vitest y Testing Library}
\label{subsec:vitest}

Las pruebas unitarias y de integración usan \textbf{Vitest 4} sobre un entorno
\texttt{jsdom}, con \textbf{Testing Library} para interactuar con los componentes como lo
haría un usuario (consultas por rol accesible, eventos de usuario). La tabla~\ref{tab:tooling}
resume la cadena de herramientas de calidad del frontend.

\vspace{0.2cm}
\noindent
\renewcommand{\arraystretch}{1.3}
\fontsize{8.5}{10.2}\selectfont\sffamily
\begin{tabularx}{\textwidth}{|L{3.4cm}|L{3.0cm}|Y|}
\hline
\rowcolor{headerTabla}
\sffamily\textbf{Herramienta} & \sffamily\textbf{Tipo} & \sffamily\textbf{Garantía que aporta} \\
\hline
TypeScript (\texttt{tsc}) & Verificación de tipos & Errores de contrato detectados en compilación. \\ \hline
\rowcolor{grisTecnico}
ESLint & Análisis estático & Reglas de \textit{hooks} y convenciones de código. \\ \hline
Vitest + Testing Library & Pruebas & Comportamiento de componentes y flujos. \\ \hline
\rowcolor{grisTecnico}
Vite (\texttt{build}) & Empaquetado & Optimización y verificación previa al despliegue. \\ \hline
\end{tabularx}
\label{tab:tooling}

\begin{NotaTecnica}
La secuencia \texttt{tsc} $\rightarrow$ \texttt{vite build} convierte la verificación de
tipos en una \textbf{barrera de despliegue}: el empaquetado monolítico del
capítulo~\ref{ch:backend} (sección~\ref{subsec:empaquetado},
pág.~\pageref{subsec:empaquetado}) solo incrusta el \textit{bundle} si el frontend compila
sin errores de tipo. Con esto se cierra el diseño detallado del frontend; el
capítulo~\ref{ch:basedatos} (pág.~\pageref{ch:basedatos}) documenta el motor de datos que
sustenta toda la plataforma.
\end{NotaTecnica}       % 5.8 Calidad del código y tooling
